\section{问题重述}

\subsection{问题背景}

大语言模型的训练可以类比为培养学生：参数量 $N$ 对应“脑容量”，训练 token 数 $D$ 对应“阅读量”，训练算力 $C$ 对应“教育经费”，数据质量 $Q$ 对应“教材质量”，领域配比 $\bp$ 对应“各科课时分配”。经典标度律 $L(N,D)=E+AN^{-\alpha}+BD^{-\beta}$ 描述验证集交叉熵损失随 $N$、$D$ 的幂律下降\cite{hoffmann,bahri}，但不含 $Q$ 与 $\bp$。高质量语料日益稀缺（“数据墙”），而 OpenScholar-8B 在若干任务上优于 GPT-4o\cite{asai}，这说明数据侧因素不可忽视。扩大规模、提升质量与延长上下文都要消耗算力，调整配比虽不增加总开销，却会改变算力的使用效果。赛题提供附件 A（质量信号与配比实验）、B（标度律观测）、C（排行榜、元数据与逐任务评测），要求在算力约束下研究资源配置。

\subsection{问题描述}

\textbf{问题一：}基于 A1–A3 的 22 个质量信号建立样本级与领域级综合质量分 $Q$，并与扩展集对照；定义并诊断指标冲突、分析主要分歧，在扩展集复验；基于 A4–A15 建立 17 域配比 $\bp$（单纯形约束）与 13 个验证域交叉熵损失的定量关系，在检验集 A6–A11 上验证，并用估算/外推表 A12–A15 讨论适用边界。

\textbf{问题二：}以问题一输出的 $Q$、$\bp$ 为接口，综合附件 B 建立含 $N,D,Q,\bp$ 的广义标度律；分析各因素的边际效用与弹性、领域间的替代或互补关系，推导质量提升与规模扩大可以相互替代的条件，并完成参数估计与检验。

\textbf{问题三：}总成本 $6ND+D[g(Q)-g(Q_0)]_++\eta NDL_{\mathrm{ctx}}\le C$ 下求最优资源配置，比较三类质量成本函数；解析给出 $L_{\mathrm{ctx}}^{\mathrm{crit}}=6/\eta$，并在 C7 的可行上下文长度上做敏感性分析；给出预算跨越量级时“结构性转移”的数学定义与识别方法。

\textbf{问题四：}基于附件 C 分离并量化规模扩张与非规模技术进步的贡献，说明能力度量、开源筛选、模型类型与时间轴四项口径；建立分层的 Loss–Benchmark 映射，在算力增长放缓情景下预测 12/24 个月开源模型能力前沿，并给出不确定性。

\section{问题分析}

\subsection{总体思路}

四问构成一条递进链（图~\ref{fig:flow}）：问题一把数据侧压缩为两个可注入的量，即连续质量分 $Q$ 与配比—损失关系；问题二把它们嵌入标度律，得到 $L(N,D,Q,\bp)$；问题三在成本约束下优化该损失；问题四用桥接映射把损失翻译为能力得分，再结合时间项做分解与预测。各问的产出按 P1–P4 接口文件组织，便于逐项核查上游输入和下游计算。

数据使用遵循\textbf{证据分层}：来源观测与公开基准（A1–A11、B1、B4、B5、C1–C4、C7、C8）支撑相应分析；其中 B1 呈公式生成特征，仅用于经典律校准。半合成与插值数据（B2、B3、B6–B8）用于形式比较、参数估计和有限留出检验，并明示局限；外推生成数据（A12–A15、B10）只用于适用边界讨论。方法选择遵循\textbf{可检验}原则：质量分的预处理与权重在 A1 拟合并冻结用于 A2/A3；质量冲突在统一参照下作评分后诊断，候选阈值对应明确的筛查预算；配比函数形式用留出数据比较，解析结果与数值结果互验，预测先做历史回测。

\begin{figure}[H]
\centering\small
\setlength{\tabcolsep}{3pt}
\begin{tabular}{|c|}\hline\makecell{\textbf{问题一}\\A1–A3 质量信号\\A4–A15 配比实验}\\\hline\end{tabular}
$\Rightarrow$
\begin{tabular}{|c|}\hline\makecell{\textbf{问题二}\\B 附件 + $Q$、$\bp$\\广义标度律}\\\hline\end{tabular}
$\Rightarrow$
\begin{tabular}{|c|}\hline\makecell{\textbf{问题三}\\成本约束优化\\结构性转移}\\\hline\end{tabular}
$\Rightarrow$
\begin{tabular}{|c|}\hline\makecell{\textbf{问题四}\\C 附件桥接\\分解与预测}\\\hline\end{tabular}\\[4pt]
{\footnotesize\begin{tabular}{llll}
接口 P1：$Q_d$、$\bp^*$、迁移矩阵 &
接口 P2：参数与 $N^*(C,Q)$ &
接口 P3：$L^*(C)$、$C_{\mathrm{crit}}$ &
接口 P4：分解与前沿
\end{tabular}}
\caption{四问递进关系与接口文件}
\label{fig:flow}
\end{figure}

\subsection{各问题的关键难点}

\textbf{问题一。}指标形态、量纲和优劣方向不同，且领域基线存在差异；冲突诊断既要识别来源间分歧，也要保留来源内异质性；配比属于成分数据，直接回归会受闭合约束影响；17 个配比域中仅 6 个有直接或近似的质量映射。对应采用 A1 拟合并冻结的域内预处理、六来源分层诊断与 P95 筛查阈值、clr 变换，以及规则特征到 A18 文本的质量桥接。

\textbf{问题二。}附件 B 各表来源与噪声差异很大，必须先审计再拟合；质量进入损失的形式（加性惩罚还是有效数据乘子）不能先验指定；附件 A 的 $Q_A$ 与附件 B 的 $Q_B$ 由不同实验和定义给出，现有损失数据不能识别二者的映射参数，需明确锚点假设并讨论敏感性。

\textbf{问题三。}约束非线性，决策变量跨越十余个数量级；最优解可能落在质量上下界的角点，“结构性转移”需要可计算的定义，而不是凭图形判断；$L_{\mathrm{ctx}}$ 是外生变量，只能取 C7 中的可行值。

\textbf{问题四。}Benchmark 与 Loss 的尺度不同，桥接数据中高可比层样本很少；规模与时间高度共线；C1 的时间窗仅 9 个月，C3 口径混杂；此外，饱和任务会掩盖前沿进展。

\section{模型假设}

\begin{enumerate}[label=\arabic*.,leftmargin=2em]
\item \textbf{质量信号可信。}附件 A 的质量信号按原始定义计算，logits 经 softmax 后具有序数意义，仅用于文档间的相对比较。
\item \textbf{小代理可迁移（待检验）。}配比对损失的排序效应在 1M–1B 参数范围内近似保持；该假设由三尺度留出集检验，检验结果决定其适用边界（\S\ref{sec:q1-scale}）。
\item \textbf{质量口径锚定。}附件 A、B 的质量映射不能由交叉熵水平直接识别。主分析在聚合可加与 B 侧质量惩罚仿射等假设下，采用 $Q_{\mathrm B}=Q_A$ 作为题面自然锚点，并声明其向 $Q_A=1$ 的延伸假设；当前 Q1 接口的等权目标质量为 $0.6609$。该锚点不是附件数据直接估计的结果，下游旧接口数值须重算。
\item \textbf{评价权重固定。}13 个验证域损失以等权汇总为目标损失，评价权重与被优化的配比严格分离；pile\_cc 单域作为副口径。
\item \textbf{成本模型。}训练、提质与注意力三项开销使用同一 $D$；$\eta=2\times10^{-4}$ 与附录 B 的成本函数参数按赛题给定。
\item \textbf{技术进步平滑。}非规模技术进步在分析窗口内可用连续时间项刻画；预测为情景条件预测，不含新架构等离散突破。
\end{enumerate}

\section{符号说明}

\begin{table}[H]
\centering\small
\caption{主要符号及含义}
\label{tab:symbol}
\begin{tabular}{p{2.6cm}p{9.6cm}l}
\toprule
符号 & 含义 & 单位 \\
\midrule
$N,\ D$ & 模型参数量、训练 token 数 & 个 \\
$C$ & 总算力预算 & FLOPs \\
$Q,\ Q_{\mathrm B}$ & 综合质量分（问题一口径）、附件 B 口径质量 & 无量纲，$[0,1]$ \\
$Q_0$ & 基线质量（问题三提质起点） & 无量纲 \\
$\bp=(p_1,\dots,p_{17})$ & 17 域训练配比，$p_i\ge0,\ \sum p_i=1$ & 无量纲 \\
$z_j(x),\ w_j$ & 文档 $x$ 第 $j$ 个标准化指标及其 CRITIC 权重 & 无量纲 \\
$B(x),\ \theta_d$ & 六来源分歧度、A1 同域 P95 筛查阈值 & 无量纲 \\
$L_v,\ L$ & 验证域 $v$ 的交叉熵损失、综合损失 & nats \\
$T_{iv}$ & 迁移矩阵：增配域 $i$ 对验证域 $v$ 损失的影响 & 无量纲 \\
$E,A,\alpha,B,\beta$ & 标度律基线参数 & — \\
$\kappa_N,\ \kappa_D,\ k_0$ & 质量—参数交互、质量—数据交互系数及质量地板系数 & $\kappa$ 无量纲，$k_0$ 为 nats \\
$h_N(Q),\ h_D(Q)$ & $1+\kappa_N(1-Q)$、$1+\kappa_D(1-Q)$，质量修正因子 & — \\
$\varepsilon_N,\varepsilon_D,\varepsilon_Q$ & 损失对 $N,D,Q$ 的弹性 & 无量纲 \\
$g(Q),\ \eta$ & 单位 token 提质成本函数、注意力开销系数 & FLOPs/token \\
$L_{\mathrm{ctx}},\ \kappa$ & 上下文长度、单位 $ND$ 成本 $\kappa=6+\eta L_{\mathrm{ctx}}$ & token，— \\
$C_{\mathrm{crit}}$ & 结构性转移的临界预算 & FLOPs \\
$S,\ \hat L$ & Benchmark 综合得分、标度律预测损失（规模状态） & 分，nats \\
$\tau_t,\ \delta$ & 每年得分技术增益、每年等效降损 & 分/年，nats/年 \\
\bottomrule
\end{tabular}
\end{table}

表中只列全文反复使用的重要符号，局部符号在首次出现处说明。

\section{模型的建立与求解}

\subsection{数据预处理}

\textbf{数据审计。}建模前逐表核对实际文件，发现五个影响建模路线的事实：(i) A10/A11 各有 64 行，按 \texttt{index} 连接；(ii) A16 的 17 域映射中只有 3 个 direct、3 个 near\_direct，其余 11 个为 inferred；(iii) B1 可由含五个待估系数的经典三项幂律拟合到 $\mathrm{RMSE}=1.5\times10^{-4}$，呈公式生成特征；(iv) 在相同 $(N,D,Q)$ 点上，B8 与 B6/B7 的损失相关系数为 $-0.03$，方向相反，因此 B8 不参与联合拟合；(v) C1 的提交日比发布日中位滞后 103.5 天，窗口仅为 2024-06 至 2025-03。

\textbf{指标展开与方向统一。}22 个原始字段展开为 25 个标量指标：分类器 logits 转为标签概率，ModernBERT 六档输出转为期望分，QuRating 拆为四个维度，DSIR 信号作按词数归一和域内长度残差化，RPS 规则统计量按字段含义区分正向、负向及适宜区间型。之后对适用的重尾变量作 $\log(1+x)$，负向指标取反，适宜区间型变量转为距 A1 同域中位数的负偏差；各域各指标按 A1 的 1\%/99\% 分位缩尾并 Min--Max 到 $[0,1]$，缺失值以 A1 同域同指标中位数填补。全部变换参数仅在 A1 拟合后\textbf{冻结}应用于 A2/A3。字段级处理见表~\ref{tab:q1_fields}。

\textbf{附件 C 预处理。}C8 共 1958 个 JSON，其中 4 个截断，其余逐任务解析为长表；C2 与 C4 按（发布日，组织）联表，保留 $t\ge2021$、$D>10^9$ 的记录；同一 C4 模型被多行命中时按 $1/$重复数加权。
